% called by main.tex
%
\chapter{Conclusiones}
\label{ch::conclusiones}

Este trabajo ha presentado una metodología de captura de datos, validación temporal y control
de costes para la predicción de corto plazo en mercados de Bitcoin de Polymarket. La
conclusión central es que una precisión direccional cercana al 90\,\% no garantiza una política
económica: con una latencia de entrada realista, el coste y la ejecución dominan el resultado.

El aprendizaje profundo aplicado a secuencias y al libro de órdenes aprende representaciones
útiles, pero no produce una política estable con unas dos semanas de datos. El especialista
tabular congelado conserva una señal modesta en cinco días intactos fuera de muestra:
$+0{,}349$ ticks por unidad al coste de referencia, 754 unidades y tres de cinco días
positivos; bajo el coste completo, el promedio es $-0{,}026$ ticks.

El análisis posterior detecta un cambio de régimen y propone un filtro de baja volatilidad.
Sobre el mismo bloque, ese subconjunto alcanza $+1{,}069$ ticks y cinco de cinco días
positivos. Su valor es el de una hipótesis prometedora, no el de una validación confirmatoria:
la regla se congela después del diagnóstico y queda pendiente de datos futuros.

Las aportaciones del trabajo son cuatro: un sistema de captura y control de calidad
multifuente; un protocolo temporal que incorpora coste y latencia medidos; evidencia
cuantitativa de la ruptura entre predicción y ejecución; y una trazabilidad que conserva tanto
los resultados negativos como la corrección de errores metodológicos. No se afirma
rentabilidad real ni se presenta un bot listo para desplegar.

\paragraph{Trabajo futuro.} La prioridad es ejecutar prospectivamente la política completa
sin retocar umbrales, con una sola regla de exposición por mercado y registro de fills. Solo
después podrá evaluarse el filtro de volatilidad. Con un horizonte de datos sustancialmente
mayor podrán reevaluarse arquitecturas como el \emph{Temporal Fusion
Transformer}~\cite{lim2021tft}. La regla que guía el proyecto permanece: la complejidad y
cualquier afirmación económica deben ganarse su sitio con evidencia verdaderamente nueva.
